% !PRINTED 332
\weberpage{332}

%<<CONT
\weq{2}{\begin{aligned}
\Delta'_\xi &= \Delta_{\xi+1} - \Delta_\xi. \\
&\cdot\ \cdot\ \cdot\ \cdot\ \cdot\ \cdot\ \cdot\ \cdot\ \cdot\ \cdot \\
\Delta_\xi^{(n-1)} &= \Delta_{\xi+1}^{(n-2)} - \Delta_\xi^{(n-2)}
\end{aligned}}

gesetzt wird,

\weq{3}{\varphi(\xi) = \varphi(0) + \Delta_0 B_1^{(\xi)} + \Delta'_0 B_2^{(\xi)} + \cdots + \Delta_0^{(n-1)} B_n^{(\xi)}.}

Der Werth der Formel liegt darin, dass die Reihe der Differenzen $\Delta_0, \Delta'_0, \Delta''_0 \ldots$ in vielen Fällen so rasch abnehmende Zahlenwerthe bietet, dass man sich mit einigen der ersten Glieder der Reihe (3) begnügen kann.

Ist nun $f(x)$ die zu untersuchende Function und $(\alpha, \beta)$ das Intervall, in dem eine der gesuchten Wurzeln liegt, so setzen wir, indem wir unter $m$ irgend eine geeignete ganze Zahl verstehen,
\[
\delta = \frac{\beta - \alpha}{m}; \qquad x = \alpha + \delta \xi,
\]
und wenn wir also
\[
\begin{aligned}
\Delta_x &= f(x + \delta) - f(x) \\
\Delta'_x &= \Delta_{x+\delta} - \Delta_x \\
&\cdot\ \cdot\ \cdot\ \cdot\ \cdot\ \cdot\ \cdot\ \cdot
\end{aligned}
\]
setzen

\weq{4}{f(x) = f(\alpha) + \frac{\Delta_\alpha}{\delta}(x - \alpha) + \frac{\Delta'_\alpha}{\delta^2}\frac{(x-\alpha)(x-\alpha-\delta)}{2} + \cdots}

Man könnte ebensogut auch von dem anderen Endpunkt $\beta$ des Intervalles ausgehen, und müsste dann in (3)
\[
x = \beta - \delta \xi
\]
setzen. Man würde dann erhalten
\[
\begin{aligned}
\Delta_x &= f(x - \delta) - f(x) \\
\Delta'_x &= \Delta_{x-\delta} - \Delta_x \\
&\cdot\ \cdot\ \cdot\ \cdot\ \cdot\ \cdot\ \cdot\ \cdot
\end{aligned}
\]

\weq{5}{f(x) = f(\beta) + \frac{\Delta_\beta}{\delta}(x - \beta) + \frac{\Delta'_\beta}{\delta^2}\frac{(x-\beta)(x-\beta+\delta)}{2} + \cdots}

Wenn man die Gleichung
\[
y = f(x)
\]
als Gleichung einer Curve in einem rechtwinkligen Coordinatensystem $x, y$ deutet, so lassen sich die Verhältnisse geometrisch veranschaulichen.

Wenn man z.\,B. die Gleichung (4) auf das Intervall von $\alpha$ bis $(\alpha + \delta)$ anwendet und sich mit der Berücksichtigung der ersten Differenz begnügt, so heisst das in der Sprache der Geometrie,
%CONT>>
